% =============================================================================
% frontmatter/abstract.tex — Sammanfattning / Abstract
% =============================================================================

\chapter*{Sammanfattning}
\addcontentsline{toc}{chapter}{Sammanfattning}

This thesis investigates how reliability, system stability, and execution performance change when a workflow orchestration system migrates from a single virtual machine using the Dagster DockerRunLauncher to a Kubernetes cluster using the K8sRunLauncher with isolated pod execution. Dagster 1.12.22 is used as a representative modern pipeline orchestration framework, and the workload is a CPU-bound, memory-intensive two-phase job. Three controlled experiments were conducted across six concurrency levels (1, 2, 3, 5, 7, and 10 concurrent jobs, three repetitions each) on local infrastructure: a 4\,vCPU / 4\,GB Multipass VM running the DockerRunLauncher and a Kind single-node cluster running the K8sRunLauncher.

The VM's job success rate degraded sharply with concurrency: 100\,\% at L1--L2, falling to 66.7\,\% at L3, 40.0\,\% at L4, 38.1\,\% at L5, and 30.0\,\% at L6. The failure mechanism was memory exhaustion: three or more concurrent containers each allocating 400\,MB exceeded the kernel's memory budget on the 4\,GB VM, triggering Linux OOM kills. Execution time for surviving jobs varied noticeably with concurrency (66.7\,s at L1--L2, rising to 92.9\,s at L3 and peaking at 181.6\,s at L4), as survivors compete for CPU and I/O with the dying containers. Kubernetes maintained 100\,\% success across all concurrency levels via per-pod cgroup limits, with consistent execution time and low variance. Pod scheduling latency was under 1\,second; container startup overhead ranged from 4.3\,s at L1 to 14.9\,s at L6 in the local Kind environment, growing with concurrency due to concurrent initialisation contention. The reliability crossover --- the first concurrency level where VM success rate drops below 95\,\% while Kubernetes maintains 100\,\% --- occurred at \textbf{L3 (3 concurrent jobs)}. At this level the VM effective delivery time (139.3\,s, accounting for reruns) already exceeds the K8s total wall-clock (69.7\,s), so both the reliability and performance crossovers align at L3. These findings give engineering teams an empirical, hardware-specific framework for migration decisions.

\vspace{1em}
\noindent\textbf{Keywords:} Kubernetes, Dagster, workflow orchestration, resource contention, pod isolation, crossover point, Kind, reliability measurement, DevOps

\clearpage
